% =============================================================================
%  Chapter 2 — Numerical Quadrature on Simplices
% =============================================================================

\chapter{Numerical Quadrature on Simplices}
\label{ch:quadrature}

Numerical integration on simplicial domains is fundamentally different from
integration on tensor-product cells (quadrilaterals, hexahedra).  On
tensor-product cells, one constructs $n^D$-point rules by taking the Cartesian
product of $n$-point Gauss--Legendre rules in each direction.  On simplices, no
such natural product structure exists; special rules must be
constructed~\cite{Stroud1971,Keast1986}.

This chapter surveys the approaches available, documents the rules used in
\texttt{fall\_n}, and motivates the Stroud conical product as the preferred
arbitrary-order method.


% ─────────────────────────────────────────────────────────────────────────
\section{Problem Statement}
\label{sec:quad_problem}

We seek to approximate the integral
\begin{equation}
  I[f] = \int_{\calS_D} f(\bxi)\,\dd\bxi
  \approx \sum_{g=1}^{n_g} w_g\, f(\bxi_g),
  \label{eq:quad_approx}
\end{equation}
where $\calS_D$ is the reference $D$-simplex~\eqref{eq:ref_simplex}, with
volume $|\calS_D| = 1/D!$.  A rule has \emph{degree of exactness}~$p$ if
$I[f]$ is computed exactly for all polynomials $f$ of total degree $\le p$.

\paragraph{Minimum requirements for finite elements.}
For a displacement-based finite element with polynomial order $k$ and a
constitutive law that is polynomial of degree $m$ in the strains:
\begin{itemize}[nosep]
  \item Stiffness matrix: needs degree $\ge 2(k-1)m$ (often $2(k-1)$ for
        linear elasticity, where $m=1$).
  \item Mass matrix (or $L^2$ projection): needs degree $\ge 2k$.
  \item For TET10 ($k=2$, $D=3$): stiffness needs degree~$\ge 2$; mass/projection
        needs degree~$\ge 4$.
\end{itemize}


% ─────────────────────────────────────────────────────────────────────────
\section{Classical Symmetric Rules}
\label{sec:classical_rules}

\subsection{Centroid Rule (Degree 1)}

The single-point centroid rule:
\begin{equation}
  \bxi_1 = \Bigl(\frac{1}{D+1}, \dots, \frac{1}{D+1}\Bigr),
  \qquad w_1 = \frac{1}{D!}\,.
  \label{eq:centroid_rule}
\end{equation}
This integrates all linear polynomials exactly.  It is adequate for constant-
strain elements (TET4) when only the stiffness matrix is needed, but
insufficient for quadratic elements.

\subsection{Hammer--Marlowe--Stroud Rules (HMS)}
\label{sec:HMS}

The Hammer--Marlowe--Stroud~\cite{HammerMarloweStroud1956} family provides low-
order symmetric rules.  The 4-point rule for the tetrahedron, used as the
default for TET10 in \texttt{fall\_n}, has:

\begin{center}
\begin{tabular}{ccl}
  \toprule
  $n_g$ & Degree & Description \\
  \midrule
  1     & 1      & Centroid \\
  4     & 2      & Vertices of a smaller inscribed tetrahedron \\
  \bottomrule
\end{tabular}
\end{center}

The 4-point degree-2 rule uses barycentric coordinates:
\begin{equation}
  \lambda = \bigl(\alpha, \beta, \beta, \beta\bigr)
  \quad\text{and permutations thereof},
  \qquad
  \alpha = \frac{5 - \sqrt{5}}{20},\;
  \beta  = \frac{5 + 3\sqrt{5}}{20}.
\end{equation}
All four weights are equal, $w_g = \frac{1}{4!} = \frac{1}{24}$, and are
\textbf{strictly positive}.

\begin{remark}
In \texttt{fall\_n}, this 4-point HMS rule is the \emph{default} quadrature
for TET10 elements (\texttt{default\_simplex\_rule<3,2>()} in
\texttt{SimplexQuadrature.hh}).  It was chosen over the 5-point Keast degree-3
rule precisely because all weights are positive\,---\,see
\cref{ch:negative_weights} for the detailed analysis.
\end{remark}


\subsection{Keast Rules}
\label{sec:keast}

Keast~\cite{Keast1986} tabulated a series of tetrahedral rules up to degree~8.
The most commonly cited are:

\begin{center}
\begin{tabular}{cccl}
  \toprule
  $n_g$ & Degree & Min weight & Comment \\
  \midrule
  1     & 1    & $+$      & Centroid \\
  4     & 2    & $+$      & Same as HMS \\
  5     & 3    & $\mathbf{-}$ & $w_1 = -\tfrac{2}{15}$ at centroid \\
  11    & 4    & $+$      & All positive \\
  14    & 5    & $+$      & All positive \\
  24    & 6    & $\mathbf{-}$ & Contains negative weights \\
  \bottomrule
\end{tabular}
\end{center}

The 5-point degree-3 rule is particularly notable:
\begin{equation}
  \text{Keast-5:}\quad
  w_1 = -\frac{2}{15}\,|\calS_3| = -\frac{2}{15} \cdot \frac{1}{6}
      = -\frac{1}{45},\qquad
  w_{2,\dots,5} = \frac{3}{40}\,|\calS_3| = \frac{1}{80}.
  \label{eq:keast5_weights}
\end{equation}
The centroid weight is \emph{negative}.  While this does not affect the
accuracy of the quadrature rule (it still integrates degree-3 polynomials
exactly), it causes severe problems in lumped $L^2$ projection on quadratic
elements, as analysed in \cref{ch:negative_weights}.


\subsection{Grundmann--M\"oller Rules}

The Grundmann--M\"oller family~\cite{GrundmannMoller1978} generates rules of
arbitrary odd degree for any simplex dimension, using a combinatorial formula.
The rule of index $s$ has degree $2s+1$ with
\[
  n_g = \binom{D + s + 1}{D + 1}
\]
points.  Unfortunately, for $s \ge 1$ (degree $\ge 3$), some weights are
\emph{guaranteed to be negative}.  This makes Grundmann--M\"oller unsuitable
as a general-purpose rule when positive weights are required.


% ─────────────────────────────────────────────────────────────────────────
\section{Stroud Conical Product Rule}
\label{sec:conical_product}

The Stroud conical product~\cite{Stroud1971,Duffy1982} transforms the simplex
integral into a product of 1D integrals via a \emph{collapsed-coordinate}
(Duffy) transformation, enabling the use of Gauss--Jacobi quadrature in each
direction.


\subsection{Duffy Transform}
\label{sec:duffy}

The transformation from \emph{collapsed coordinates}
$(t_1, \dots, t_D) \in [0,1]^D$ to simplex coordinates
$(\xi_1, \dots, \xi_D) \in \calS_D$ is defined recursively:
\begin{equation}
  \boxed{
  \begin{aligned}
    \xi_1 &= t_1, \\
    \xi_k &= t_k \prod_{j=1}^{k-1}(1 - t_j), \qquad k = 2, \dots, D.
  \end{aligned}}
  \label{eq:duffy_transform}
\end{equation}

The Jacobian of this transformation is
\begin{equation}
  \left|\frac{\partial\bxi}{\partial\mathbf{t}}\right|
  = \prod_{k=1}^{D-1} (1 - t_k)^{D-k}.
  \label{eq:duffy_jacobian}
\end{equation}

Therefore, the integral over the simplex becomes
\begin{equation}
  \int_{\calS_D} f(\bxi)\,\dd\bxi
  = \int_0^1 \cdots \int_0^1
    f\bigl(\bxi(\mathbf{t})\bigr)
    \prod_{k=1}^{D-1} (1 - t_k)^{D-k}
    \;\dd t_1 \cdots \dd t_D.
  \label{eq:duffy_integral}
\end{equation}


\subsection{Product Quadrature with Gauss--Jacobi Nodes}

The key insight is that each factor $(1 - t_k)^{D-k}$ depends only on $t_k$,
so it can be \emph{absorbed into} the 1D quadrature rule for direction $k$.
Specifically, for direction $k$ ($0$-indexed), we need a 1D quadrature rule on
$[0,1]$ with weight function $(1-t)^{\beta_k}$, where
\begin{equation}
  \beta_k = D - 1 - k, \qquad k = 0, 1, \dots, D-1.
  \label{eq:beta_definition}
\end{equation}

\begin{itemize}[nosep]
  \item $k = 0$: $\beta = D-1$ (maximal weight for the ``lowest'' direction).
  \item $k = D-2$: $\beta = 1$ (linear weight).
  \item $k = D-1$: $\beta = 0$ (plain Gauss--Legendre on $[0,1]$).
\end{itemize}

The integral~\eqref{eq:duffy_integral} becomes
\begin{equation}
  \int_{\calS_D} f(\bxi)\,\dd\bxi
  = \prod_{k=0}^{D-1} \left(
    \int_0^1 g_k(t_k)\,(1 - t_k)^{\beta_k}\,\dd t_k
  \right),
  \label{eq:product_integral}
\end{equation}
where $g_k$ encodes the dependence of $f$ on $t_k$ after integrating all other
variables.  Applying $n$-point Gauss--Jacobi quadrature on $[0,1]$ with weight
$(1-t)^{\beta_k}$ to each direction gives a product rule with:
\begin{equation}
  \boxed{
  n_g = n^D \text{ points}, \qquad
  \text{degree of exactness} = 2n - 1.}
  \label{eq:conical_cost}
\end{equation}

\begin{proposition}[All weights positive]
\label{prop:positive_weights}
Since each 1D Gauss--Jacobi rule has strictly positive weights (the weight
function $(1-t)^{\beta}$ is positive on $(0,1)$ for $\beta > -1$, and all
Gauss--Jacobi weights are positive for any $\alpha, \beta > -1$), the product
rule inherits positivity:
\begin{equation}
  w_g = \prod_{k=0}^{D-1} w_{k,i_k} > 0 \qquad \forall\, g.
  \label{eq:positive_product_weight}
\end{equation}
This is the fundamental advantage of the conical product over symmetric
simplex rules.
\end{proposition}


\subsection{Cost Comparison}

\begin{center}
\begin{tabular}{lccl}
  \toprule
  Rule & Points & Degree & Weights \\
  \midrule
  HMS 4-pt (TET)          & 4    & 2 & All $+$ \\
  Keast 5-pt              & 5    & 3 & \textbf{1 negative} \\
  Keast 11-pt             & 11   & 4 & All $+$ \\
  Conical $n\!=\!2$, $D\!=\!3$ & 8   & 3 & All $+$ \\
  Conical $n\!=\!3$, $D\!=\!3$ & 27  & 5 & All $+$ \\
  Conical $n\!=\!4$, $D\!=\!3$ & 64  & 7 & All $+$ \\
  \bottomrule
\end{tabular}
\end{center}

The conical product uses more points than optimal symmetric rules for the same
degree, but it guarantees positivity and extends to \emph{any} polynomial
degree.  For the TET10 stiffness matrix (degree $\ge 2$), the 4-point HMS rule
is sufficient and more economical; the conical product is available as a
drop-in replacement when higher accuracy is needed.


% ─────────────────────────────────────────────────────────────────────────
\section{Gauss--Jacobi Quadrature via the Golub--Welsch Algorithm}
\label{sec:golub_welsch}

The 1D Gauss--Jacobi rules needed for the conical product are computed at
library initialisation time using the Golub--Welsch
algorithm~\cite{GolubWelsch1969}.

\subsection{Jacobi Polynomials}
\label{sec:jacobi_poly}

The Jacobi polynomials $P_n^{(\alpha,\beta)}(x)$ are orthogonal on $[-1,1]$
with respect to the weight function
\begin{equation}
  w(x) = (1 - x)^\alpha (1 + x)^\beta, \qquad \alpha, \beta > -1.
  \label{eq:jacobi_weight}
\end{equation}
They satisfy the three-term recurrence~\cite{Szego1975,Gautschi2004}:
\begin{equation}
  P_{n+1}^{(\alpha,\beta)}(x)
  = (A_n\, x + B_n)\, P_n^{(\alpha,\beta)}(x)
  - C_n\, P_{n-1}^{(\alpha,\beta)}(x),
  \label{eq:3term_recurrence}
\end{equation}
where the coefficients $A_n$, $B_n$, $C_n$ are explicit functions of $n$,
$\alpha$, $\beta$.

The \emph{monic} Jacobi polynomials $\hat{p}_n$ satisfy
\[
  \hat{p}_{n+1}(x) = (x - a_n)\,\hat{p}_n(x) - b_n\,\hat{p}_{n-1}(x),
\]
with recurrence coefficients~\cite{Gautschi2004}:
\begin{align}
  a_n &= \frac{\beta^2 - \alpha^2}{(2n+\alpha+\beta)(2n+\alpha+\beta+2)},
  \label{eq:an_coeff} \\
  b_n &= \frac{4n(n+\alpha)(n+\beta)(n+\alpha+\beta)}
              {(2n+\alpha+\beta)^2(2n+\alpha+\beta+1)(2n+\alpha+\beta-1)}.
  \label{eq:bn_coeff}
\end{align}


\subsection{The Golub--Welsch Algorithm}
\label{sec:gw_algorithm}

The Golub--Welsch algorithm~\cite{GolubWelsch1969} computes the $n$-point
Gauss quadrature rule from the eigendecomposition of the \emph{symmetric
tridiagonal Jacobi matrix} $\mathbf{T} \in \R^{n \times n}$:
\begin{equation}
  \mathbf{T} = \begin{pmatrix}
    a_0       & \sqrt{b_1} &            &        \\
    \sqrt{b_1} & a_1       & \sqrt{b_2} &        \\
                & \ddots    & \ddots     & \ddots \\
                &           & \sqrt{b_{n-1}} & a_{n-1}
  \end{pmatrix}.
  \label{eq:jacobi_matrix}
\end{equation}

\begin{itemize}[nosep]
  \item The \textbf{eigenvalues} $\lambda_1, \dots, \lambda_n$ of
        $\mathbf{T}$ are the \emph{quadrature nodes} (zeros of the degree-$n$
        monic Jacobi polynomial).
  \item The \textbf{weights} are $w_i = \mu_0\, v_{1,i}^2$, where
        $v_{1,i}$ is the first component of the $i$-th normalised eigenvector
        and $\mu_0$ is the zeroth moment of the weight function:
        \begin{equation}
          \mu_0 = \int_{-1}^{1} (1-x)^\alpha (1+x)^\beta\,\dd x
               = \frac{2^{\alpha+\beta+1}\,\Gamma(\alpha+1)\,\Gamma(\beta+1)}
                      {\Gamma(\alpha+\beta+2)}.
          \label{eq:mu0}
        \end{equation}
\end{itemize}

\begin{remark}[Why Golub--Welsch instead of Newton iteration]
An earlier implementation attempted to find the Jacobi polynomial zeros via a
Newton iteration with Tricomi-type initial guesses.  This approach failed for
the $\alpha = D-1$ cases needed by the conical product (e.g., $\alpha = 2$ for
$D = 3$, direction $k=0$): the initial guesses were too inaccurate, and the
iteration converged to \emph{duplicate roots} instead of distinct nodes.

The Golub--Welsch algorithm is unconditionally robust: it requires no initial
guesses and produces all nodes and weights simultaneously from a single
eigenvalue decomposition.  It is implemented using Eigen's
\texttt{SelfAdjointEigenSolver}, which is $O(n^2)$ for the tridiagonal case.
\end{remark}


\subsection{Mapping to $[0,1]$}

The Gauss--Jacobi rule on $[-1,1]$ with weight $(1-x)^\alpha$ is mapped to
$[0,1]$ via $x = 2t - 1$:
\begin{equation}
  \int_0^1 f(t)\,(1-t)^\beta\,\dd t
  = \frac{1}{2^{\beta+1}} \int_{-1}^{1} f\!\left(\frac{1+x}{2}\right)
    (1-x)^\beta\,\dd x.
  \label{eq:map_to_01}
\end{equation}
The nodes and weights transform as
\begin{equation}
  t_i = \frac{1 + x_i}{2}, \qquad
  w_i^{[0,1]} = \frac{w_i^{[-1,1]}}{2^{\beta+1}}.
  \label{eq:weight_map}
\end{equation}
This is implemented in \texttt{gauss\_jacobi\_01()} in
\texttt{StroudConicalProduct.hh}.


% ─────────────────────────────────────────────────────────────────────────
\section{Implementation: \texttt{StroudConicalProduct.hh}}
\label{sec:conical_impl}

The implementation resides in
\texttt{src/numerics/numerical\_integration/StroudConicalProduct.hh} and
consists of three layers:

\begin{enumerate}[leftmargin=2em]
  \item \textbf{\texttt{gauss\_jacobi(n, $\alpha$, $\beta$)}} ---
        Returns an $n$-point Gauss--Jacobi rule on $[-1,1]$ via Golub--Welsch.
        Uses Eigen's \texttt{SelfAdjointEigenSolver} on the tridiagonal Jacobi
        matrix.

  \item \textbf{\texttt{stroud\_conical\_product<D>(n\_per\_dir)}} ---
        Builds $D$ one-dimensional Gauss--Jacobi rules (with
        $\beta_k = D-1-k$), takes their tensor product in collapsed
        coordinates, and maps to simplex coordinates via the Duffy
        transform~\eqref{eq:duffy_transform}.

  \item \textbf{\texttt{ConicalProductIntegrator<TopDim, NPerDir>}} ---
        Compile-time-sized integrator class.  The rule is computed once at
        program startup (static \texttt{inline const} initialiser) and stored
        in a fixed-size \texttt{std::array}.  Exposes the same interface as
        all other integrators in \texttt{fall\_n}:
        \begin{itemize}[nosep]
          \item \texttt{static constexpr num\_integration\_points = NPerDir$^{\text{TopDim}}$}
          \item \texttt{reference\_integration\_point(i)} returns a $\texttt{std::span<const double>}$
          \item \texttt{weight(i)} returns $w_i$
          \item \texttt{operator()(f)} evaluates $\sum w_i f(\bxi_i)$
        \end{itemize}
\end{enumerate}


% ─────────────────────────────────────────────────────────────────────────
\section{Default Quadrature Choices in \texttt{fall\_n}}
\label{sec:default_rules}

\begin{center}
\begin{tabular}{llccl}
  \toprule
  Element & Rule & Points & Degree & Source \\
  \midrule
  TET4  & HMS 1-pt  & 1  & 1 & \texttt{default\_simplex\_rule<3,1>()} \\
  TET10 & HMS 4-pt  & 4  & 2 & \texttt{default\_simplex\_rule<3,2>()} \\
  TRI3  & HMS 1-pt  & 1  & 1 & \texttt{default\_simplex\_rule<2,1>()} \\
  TRI6  & HMS 3-pt  & 3  & 2 & \texttt{default\_simplex\_rule<2,2>()} \\
  HEX27 & Gauss $3^3$ & 27 & 5 & Tensor-product Gauss--Legendre \\
  \bottomrule
\end{tabular}
\end{center}

The conical product integrators (\texttt{ConicalProductIntegrator<3,2>},
\texttt{ConicalProductIntegrator<3,3>}, etc.) are available as drop-in
replacements when higher-degree or higher-accuracy integration is needed,
with the guarantee that all weights remain positive.
